\section{Plan Mode：先想清楚再动手}

\begin{knowledgebox}{Plan Mode 是什么？}
在 Plan Mode 下，Claude Code 只能读取文件和搜索代码，不能做任何修改。它会探索代码库、分析需求，然后给出一个详细的实现计划。进入方式包括：在提示词中说明"先给我一个计划"、使用 \texttt{Shift+Tab} 切换模式，或在 \texttt{CLAUDE.md} 中配置默认行为。
\end{knowledgebox}

\subsection{Plan Mode 的正确工作流}

\begin{importantbox}{官方推荐四步流程：探索→规划→实现→验证}
\begin{enumerate}[leftmargin=1.5em]
  \item 让 Claude 生成初始计划
  \item 按 \texttt{Ctrl+G} 在编辑器中打开并编辑计划——删掉不同意的部分，补充自己的想法
  \item 切回正常模式（\texttt{Shift+Tab}），让 Claude 按修改后的计划执行
  \item 执行完毕后验证结果（跑测试、检查输出）
\end{enumerate}
\textbf{关键认知}：改一个计划只需要一句话，改已经写好的代码要花十倍的时间。
\end{importantbox}

\subsection{何时不需要 Plan Mode}

判断标准很简单：如果任务的实现方式只有一种（改变量名、修 typo、加一行日志），直接做。如果有多种选择，先规划。

\subsection{本章小结}

Plan Mode 是复杂任务的必经之路。先规划再动手，可以大幅减少返工。核心操作是 \texttt{Ctrl+G} 编辑计划。

%% ====================================================================
